\sscp{Postpositions in prepositional SVO languages}{02-03-01-03}
Other structures are also found sporadically in the region that
are associated cross-linguistically with verb-final OV typologies
such as postpositions \citep{Dryer2013a,Dryer2013f},
and are not typical of Austronesian languages in western Indonesia or the Philippines.
Postpositions in the target region are not widespread,
but they are found in languages that are mainly prepositional.
For example, \ili{Dela} (western Rote) is predominantly prepositional,
but verbs involving motion through space or time can also trigger
one of three postpositional verbs indicating direction of motion
\citep[291--293]{Tamelan2021} as in examples \qf{02-29}--\qf{02-31}.
\ili{As} is typical with predominantly head-marking languages,
many \tsc{Timor-Babar} languages (which includes languages of
Rote) inflect most adpositions
and complementisers, as well as verbs.\footnote{
		In some of these
		languages many adpositions and complementisers are morphologically indistinguishable from
		verbs and the basis for distinguishing all three word classes is not always clear.
		For instance, \citet{Edwards2020} treats all these kinds
		of words in \ili{Kotos} \ili{Amarasi} as verbs.}
It is assumed that the
postpositions developed through serial verb constructions,
but they nevertheless appear \it{after} the adpositional complement,
in a slot that is unusual for other serial verb constructions
in the Rote languages.

\newpage
\begin{exe}
	\ex{\ili{Dela} [row]\label{ex:02-29}}\jambox*{\citep{UBB2020a}}\vspace{-1mm}
	\sn{\gll	ɓasa ma, ana lao na-kandoo seriʔ rulu \tbf{n-eu}\\
						after and \tsc{3sg.nom} go \tsc{3sg}-continue side \ili{east} \tsc{3sg}-go\\
			\glt	`After that, he continued journeying to the \ili{east}.'
						}\label{ex:}
	\ex{\gll	ɓoɓon ma, ana ɓaliʔ mia osi \tbf{n-ema}\\
						after and \tsc{3sg.nom} return from garden \tsc{3sg}-come\\
			\glt	`That afternoon, he came home from the field.'
						}\label{ex:02-30}
	\ex{\gll	huu naa, au fee ne \tbf{n-eti}\\
						case \tsc{dist} \tsc{1sg} give \tsc{3sg.acc} \tsc{3sg}-toward.addressee\\
			\glt	`Therefore, I sent him to you.'
						}\label{ex:02-31}
\end{exe}

Postpositions, while not prevalent in the Austronesian languages
in the region, are found in both \tsc{Rote-Meto}
languages as well as \tsc{Flores-Lembata}
languages (see \citeauthor{Fricke2019} [\citeyear[44ff]{Fricke2019}], where they
are analysed as clause-final deictic motion verbs).
While many other \tsc{Timor-Babar} languages also have extensive
use of serial verbs
(\citealp{JacobGrimes2011}, \citealp[311--319]{Edwards2020}, \citealp[413--439]{Tamelan2021}),
many of those do not have postpositions,
so having serial verbs alone is not a sufficient explanation
for postpositions, although it must certainly be part of the discussion.
Once again, contact with typologically SOV Papuan languages for
which postpositions are the norm \citep{Klamer2017} provides
the most adequate explanation for this typological mix of languages
in a limited region that have mostly prepositions
and also a few postpositions.